% chapters/appendix_b.tex
\chapter{Implementation Details}
\label{app:implementation}

% ============================================================
% APPENDIX OVERVIEW
% ============================================================
% This appendix contains:
% - Code snippets for key analyses
% - Software configuration details
% - Data processing pipeline


% ============================================================
% B.1 SOFTWARE ENVIRONMENT
% ============================================================
\section{Software Environment}
\label{app:software}

% CONTENT:
%
% TABLE: Software versions
% Package | Version | Purpose
% Python | 3.X | Base language
% PyTorch | X.X.X | Deep learning framework
% MNE | X.X.X | EEG analysis
% Captum | X.X.X | Attribution methods
% NumPy | X.X.X | Numerical computing
% SciPy | X.X.X | Scientific computing
% Matplotlib | X.X.X | Visualization
% etc.
%
% Requirements.txt or environment.yml equivalent


% ============================================================
% B.2 ERDS COMPUTATION
% ============================================================
\section{ERDS Computation Code}
\label{app:erds_code}

% CONTENT:
%
% Key code snippets for ERDS analysis
% Use listings package for syntax highlighting

\begin{lstlisting}[language=Python, caption={ERDS computation using MNE}, label={lst:erds}]
import mne
import numpy as np

def compute_erds(epochs, freqs, baseline, tmin, tmax):
    """
    Compute Event-Related Desynchronization/Synchronization.
    
    Parameters
    ----------
    epochs : mne.Epochs
        Epoched EEG data
    freqs : array-like
        Frequencies of interest
    baseline : tuple
        Baseline period (start, end) in seconds
    tmin, tmax : float
        Time window for analysis
    
    Returns
    -------
    erds : array
        ERDS values (channels x frequencies x times)
    """
    # Time-frequency decomposition
    power = mne.time_frequency.tfr_morlet(
        epochs, 
        freqs=freqs,
        n_cycles=freqs / 2,  # Adaptive cycles
        return_itc=False,
        average=True
    )
    
    # Extract baseline power
    baseline_mask = (power.times >= baseline[0]) & \
                    (power.times <= baseline[1])
    baseline_power = power.data[..., baseline_mask].mean(axis=-1, 
                                                         keepdims=True)
    
    # Compute ERDS percentage
    erds = (baseline_power - power.data) / baseline_power * 100
    
    return erds, power.times, freqs
\end{lstlisting}


% ============================================================
% B.3 FILTER EXTRACTION AND VISUALIZATION
% ============================================================
\section{Filter Analysis Code}
\label{app:filter_code}

\begin{lstlisting}[language=Python, caption={Temporal filter frequency response}, label={lst:temporal_filter}]
import numpy as np
import matplotlib.pyplot as plt

def analyze_temporal_filters(model, fs=250):
    """
    Extract and analyze temporal filter frequency responses.
    
    Parameters
    ----------
    model : nn.Module
        Trained EEGNet model
    fs : float
        Sampling frequency in Hz
    
    Returns
    -------
    freqs : array
        Frequency axis
    responses : array
        Magnitude responses (n_filters x n_freqs)
    """
    # Extract temporal conv weights
    # Adjust layer name based on your implementation
    weights = model.temporal_conv.weight.detach().cpu().numpy()
    # Shape: (F1, 1, 1, kernel_length)
    
    n_filters = weights.shape[0]
    kernel_length = weights.shape[-1]
    
    # Zero-pad for frequency resolution
    n_fft = 512
    
    responses = []
    for i in range(n_filters):
        kernel = weights[i, 0, 0, :]
        
        # Compute frequency response
        freq_response = np.abs(np.fft.rfft(kernel, n=n_fft))
        responses.append(freq_response)
    
    responses = np.array(responses)
    freqs = np.fft.rfftfreq(n_fft, d=1/fs)
    
    return freqs, responses


def plot_filter_responses(freqs, responses, fs=250):
    """Plot frequency responses with band markers."""
    fig, ax = plt.subplots(figsize=(10, 6))
    
    for i, resp in enumerate(responses):
        ax.plot(freqs, resp / resp.max(), label=f'Filter {i+1}')
    
    # Mark frequency bands
    ax.axvspan(8, 13, alpha=0.2, color='green', label='Mu (8-13 Hz)')
    ax.axvspan(13, 30, alpha=0.2, color='blue', label='Beta (13-30 Hz)')
    
    ax.set_xlabel('Frequency (Hz)')
    ax.set_ylabel('Normalized Magnitude')
    ax.set_xlim(0, 40)
    ax.legend()
    
    return fig
\end{lstlisting}


\begin{lstlisting}[language=Python, caption={Spatial filter topography}, label={lst:spatial_filter}]
import mne

def plot_spatial_filters(model, info, n_filters=16):
    """
    Visualize spatial filter weights as topographic maps.
    
    Parameters
    ----------
    model : nn.Module
        Trained EEGNet model
    info : mne.Info
        MNE info object with electrode positions
    n_filters : int
        Number of spatial filters
    
    Returns
    -------
    fig : Figure
        Matplotlib figure with topographic maps
    """
    # Extract spatial conv weights
    weights = model.spatial_conv.weight.detach().cpu().numpy()
    # Shape: (F1*D, 1, n_channels, 1)
    
    n_channels = weights.shape[2]
    
    # Create grid of topoplots
    n_rows = 4
    n_cols = 4
    fig, axes = plt.subplots(n_rows, n_cols, figsize=(12, 12))
    
    for i, ax in enumerate(axes.flat):
        if i < n_filters:
            filter_weights = weights[i, 0, :, 0]
            
            mne.viz.plot_topomap(
                filter_weights, 
                info, 
                axes=ax,
                show=False,
                contours=0
            )
            ax.set_title(f'Filter {i+1}')
        else:
            ax.axis('off')
    
    plt.tight_layout()
    return fig
\end{lstlisting}


% ============================================================
% B.4 INTEGRATED GRADIENTS IMPLEMENTATION
% ============================================================
\section{Integrated Gradients Code}
\label{app:ig_code}

\begin{lstlisting}[language=Python, caption={Integrated Gradients analysis}, label={lst:ig}]
import torch
from captum.attr import IntegratedGradients
import numpy as np

def compute_attributions(model, data_loader, device='cuda'):
    """
    Compute Integrated Gradients attributions for all samples.
    
    Parameters
    ----------
    model : nn.Module
        Trained model (must be in eval mode)
    data_loader : DataLoader
        DataLoader with test samples
    device : str
        Device to use
    
    Returns
    -------
    attributions : dict
        Attributions organized by class
    """
    model.eval()
    model.to(device)
    
    ig = IntegratedGradients(model)
    
    attributions = {'left': [], 'right': []}
    
    for batch_x, batch_y in data_loader:
        batch_x = batch_x.to(device)
        batch_x.requires_grad = True
        
        # Create baseline (zeros)
        baseline = torch.zeros_like(batch_x)
        
        for i in range(len(batch_x)):
            x = batch_x[i:i+1]
            y = batch_y[i].item()
            
            # Get prediction
            with torch.no_grad():
                pred = model(x).argmax(dim=1).item()
            
            # Only analyze correctly classified samples
            if pred == y:
                # Compute attribution
                attr = ig.attribute(
                    x,
                    baselines=baseline[i:i+1],
                    target=pred,
                    n_steps=100
                )
                
                attr_np = attr.detach().cpu().numpy()
                
                # Store by class
                class_name = 'left' if y == 0 else 'right'
                attributions[class_name].append(attr_np)
    
    # Convert to arrays
    for key in attributions:
        attributions[key] = np.concatenate(attributions[key], axis=0)
    
    return attributions


def aggregate_attributions(attributions, method='mean'):
    """
    Aggregate attributions across trials.
    
    Parameters
    ----------
    attributions : dict
        Per-trial attributions from compute_attributions
    method : str
        'mean' or 'abs_mean'
    
    Returns
    -------
    aggregated : dict
        Aggregated attributions per class
    """
    aggregated = {}
    
    for class_name, attrs in attributions.items():
        if method == 'mean':
            aggregated[class_name] = attrs.mean(axis=0)
        elif method == 'abs_mean':
            aggregated[class_name] = np.abs(attrs).mean(axis=0)
    
    return aggregated
\end{lstlisting}


% ============================================================
% B.5 ELECTRODE CONFIGURATION
% ============================================================
\section{Electrode Configuration}
\label{app:electrodes}

% CONTENT:
%
% TABLE: Full electrode list
% Electrode | Position (10-10) | Region | ROI Assignment
%
% E.g.:
% Fp1 | Frontal polar left | Frontal | Artifact ROI
% Fp2 | Frontal polar right | Frontal | Artifact ROI
% F3 | Frontal left | Frontal | Artifact ROI
% ...
% C3 | Central left | Sensorimotor | Sensorimotor ROI
% C4 | Central right | Sensorimotor | Sensorimotor ROI
% ...
% O1 | Occipital left | Occipital | Visual ROI
% O2 | Occipital right | Occipital | Visual ROI
%
% FIGURE: Electrode montage diagram


% ============================================================
% B.6 HYPERPARAMETERS
% ============================================================
\section{Model and Analysis Hyperparameters}
\label{app:hyperparameters}

% CONTENT:
%
% TABLE: EEGNet hyperparameters
% Parameter | Value | Description
% F1 | 8 | Number of temporal filters
% D | 2 | Depth multiplier
% F2 | 16 | Number of separable filters
% kernel_length | 64 | Temporal kernel size (samples)
% dropout | 0.3 | Dropout rate
% ...
%
% TABLE: Analysis hyperparameters
% Parameter | Value | Description
% IG n_steps | 100 | Integration steps
% Wavelet n_cycles | freqs/2 | Wavelet cycles
% ...
